\subsection{21.115: Fourier support and class-two linearization}
\label{prior:21.115}

Sambale's preprint of 8 September 2026, two days before the
experiment, proves the general complement bound for unions
of cosets, including mixed left and right cosets
\cite{sambale-cosets}. The experiment first developed the
following restricted arguments, then found that resolution
and excluded the problem from its candidate count.
The earlier Sambale--T\u arn\u auceanu paper treated several
special cases \cite{sambale-tarnauceanu}; the character-product
method below also appears in Szegedy's work \cite{szegedy}.

If cosets \(g_iH_i\), \(1\le i\le n\), do not cover a finite
abelian group \(A\), choose \(x_0\) outside their union.
Characters separate points in \(A/H_i\), so choose \(\chi_i\)
trivial on \(H_i\) with \(\chi_i(x_0)\ne\chi_i(g_i)\).
The nonzero function
\[
 f(x)=\prod_{i=1}^n(\chi_i(x)-\chi_i(g_i))
\]
vanishes on the union and is a linear combination of at most
\(2^n\) distinct characters. For a nonzero combination
\(f=\sum_{\chi\in T}a_\chi\chi\), orthogonality and
Cauchy--Schwarz give
\[
 |A|\sum|a_\chi|^2=\|f\|_2^2
 \le|\operatorname{supp}f|\|f\|_\infty^2
 \le|\operatorname{supp}f|\,|T|\sum|a_\chi|^2.
\]
Thus the complement has at least \(|A|/2^n\) elements.

The same deduction applies to a class-two finite group
whose derived subgroup has odd order. On its underlying
set define \(x+y=xy[x,y]^{-1/2}\), taking the unique half
only in the central odd-order derived subgroup.
Class-two commutator identities show this is abelian:
both triple sums equal
\(xyz[x,y]^{-1/2}[x,z]^{-1/2}[y,z]^{-1/2}\).
Every original subgroup is additive, since the half of
its commutator is a power of that commutator.
For fixed \(g\), \(T_g(x)=[g,x]\) is an additive central-valued
map with \(T_g^2=0\); hence \(L_g=1+\frac12T_g\) is an
additive automorphism, and \(gx=g+L_g(x)\).
Every original coset is therefore an affine coset in the
same abelian group. Applying the Fourier argument proves
the bound. This avoids any assumption that the whole
group has odd order, but does not cover every class-two
2-group. The general theorem is credited to the earlier
preprint, not inferred from these restricted cases.
See \artifact{research/21.115-partial-proof.md}.

\subsection{16.60: a tensor proof of an existing character bound}
\label{prior:16.60}

The full affirmative answer was stated in the May 2026
preprint of Yerrapati--Dixit--Shukla
\cite[Theorem~3.1]{yerrapati-dixit-shukla}.
For a finite group \(G\) and any automorphism \(\alpha\),
the number of elements inverted by \(\alpha\) is at most
\(T(G)=\sum_{\chi\in\operatorname{Irr}(G)}\chi(1)\).
The following proof supplies the needed estimate without
treating a character value on a product as a product of
character values.

For a unitary irreducible representation \(\rho\) on \(V\),
put
\[
 E=\frac1{|G|}\sum_{g\in G}\rho(g)\otimes\rho(\alpha(g)),
 \qquad F(v\otimes w)=w\otimes v.
\]
The operator \(E\) is the orthogonal projection onto invariants
of \(V\otimes(V\circ\alpha)\). Its rank is
\(\dim\operatorname{Hom}_G(V^*,V\circ\alpha)\), at most one
by Schur's lemma. The trace identity
\(\tr((A\otimes B)F)=\tr(AB)\) gives
\[
 \iota_\alpha(\chi)
 =\frac1{|G|}\sum_g\chi(g\alpha(g))=\tr(EF),\qquad
 |\iota_\alpha(\chi)|\le\operatorname{rank}E\le1.
\]
The inequality uses an orthonormal basis of \(\operatorname{im}E\)
and unitarity of \(F\); \(E\) need not commute with \(F\).
The regular-character identity now yields
\[
 |\{g:\alpha(g)=g^{-1}\}|
   =\sum_\chi\chi(1)\iota_\alpha(\chi)\le T(G).
\]
No involutivity of \(\alpha\) is needed.
Alternatively, every inverted element lies in
\(\operatorname{Fix}(\alpha^2)\), where the restriction of
\(\alpha\) is involutive. The classical twisted indicator
bound there and \(T(H)\le T(G)\) for \(H\le G\) imply the
same result; the latter inequality follows by restricting
the direct sum of all irreducible \(G\)-modules and using
Frobenius reciprocity.
The source audit identifies an unexplained scalar
Cauchy--Schwarz passage in its Lemma~2.19; this tensor
argument establishes the needed statement independently of
that passage, without assigning priority to the answer.
See \artifact{research/16.60-known-consequence.md}.
